\chapter{总结与展望}
\label{chapter:conclusion}

本文提出了一种全新的隐私保护程序验证框架，实现了从程序实现（SSA）到 SAT 公式生成，再到最终属性验证（SAT/UNSAT）的全流程隐私保护。该框架通过两阶段协议设计，有效解决了验证透明性与知识产权保护之间的根本矛盾。

\section{主要贡献总结}
\begin{enumerate}
    \item \textbf{高效的转换协议}：在阶段一，利用 HyperPlonk 优化的 zkWASM 实现了 SSA 程序到 SAT 公式的快速转换。实验证明，该方案在处理大规模程序时仍能保持线性扩展性，性能提升约 3 倍，显著降低了系统的启动延迟。
    \item \textbf{宽度无关的 UNSAT 验证}：在阶段二，提出了挑战值编码技术，将变长子句压缩为单一域元素。实验数据表明，验证时间不再随子句宽度增加而增长，彻底解决了宽子句带来的性能瓶颈。
    \item \textbf{常数级的 SAT 验证}：设计了双矩阵赋值结构，利用并行约束使得 SAT 验证复杂度独立于公式规模。无论公式多复杂，验证开销始终保持在常数级别，具备极高的可扩展性。
\end{enumerate}

\section{未来展望}
未来的工作将集中在以下几个方向：首先，进一步优化大规模并发程序的处理能力，探索使用递归证明（Recursive Proofs）来进一步压缩证明生成时间；其次，扩展框架以支持浮点数运算和更复杂的算术理论，以适应更广泛的程序验证需求；最后，研究基于硬件加速（FPGA/GPU）的证明生成方案，推动该技术在实际工业场景中的落地应用。